\subsection{经血逆流与子宫内膜异位症：经期性交会加重吗？}

几乎所有关于经期性行为的讨论，最后都会绕到这个问题上："经血逆流会造成子宫内膜异位症，所以经期别做。"这个说法一半来自真实的医学理论，一半是被误用的推论。这一节把三层分开讲——理论是什么、理论支持什么、理论不能推出什么。

\textbf{一、逆流理论是什么}

\begin{itemize}
\item \textbf{Sampson 学说（1927 年）}：月经期部分经血连同脱落的子宫内膜细胞，经输卵管向腹腔方向逆流，种植在盆腔腹膜、卵巢等部位，形成异位病灶。它至今仍是子宫内膜异位症\textbf{主流的发病学说之一}；
\item \textbf{支持它的证据}：经期可在腹腔液中检出内膜细胞；异位病灶的分布与逆流路径相符；与输卵管通畅情况相关的观察结果；
\item 其他并存的学说还有体腔上皮化生、免疫异常、遗传易感、干细胞来源等。也就是说，现代医学对内异症的认识是\textbf{多因素}的，而不是单一机制。
\end{itemize}

\textbf{二、这个理论不能直接推出什么}

关键在于一个经常被跳过的反证：

\begin{itemize}
\item \textbf{经血逆流在几乎所有女性身上都存在}——包括月经完全规律、从未有盆腔问题的女性，也包括因其他原因做腹腔镜手术时被偶然观察到的女性；
\item 而子宫内膜异位症发生在约十分之一左右的育龄女性身上；
\item 如果逆流是唯一原因，那么患病率应该与逆流的发生率相当。二者之间的巨大落差说明：\textbf{逆流是常见现象，而发病还需要其他条件}——例如免疫清除功能异常、遗传易感、局部微环境因素等；
\item 因此正确的表述是："逆流是发病机制的重要一环"，而\textbf{不是}"有逆流就会得内异症"。
\end{itemize}

\textbf{三、经期性交会不会导致或加重内异症}

\begin{itemize}
\item \textbf{目前的结论}：\textbf{没有证据}表明经期性交会导致子宫内膜异位症，也没有证据表明它会使已有的内异症加重；
\item \textbf{理由}：性行为过程中的机械与压力作用，与月经期自然发生的逆流量相比并不突出；同时，逆流本身既然普遍存在，就无法用"是否发生逆流"来解释少数人的发病；
\item \textbf{需要诚实说明的局限}：直接针对"经期性交与内异症关系"的研究很少（这类研究在设计上难度大、也涉及伦理考量）。所以准确的结论是"目前没有证据支持这一禁忌"，而\textbf{不是}已证明"绝对无任何影响"。医学上这两句话的差别很重要；
\item \textbf{但有明确不适的人群需要区别对待}：已确诊子宫内膜异位症或子宫腺肌症者，若在经期性交后出现明显疼痛加重，应依据自身症状调整——这不是因为逆流理论，而是因为疼痛本身就是需要尊重的信号。
\end{itemize}

\textbf{四、比"逆流"更值得排查的几件事}

经期性交后若出现下列情况，更需要排查的是具体问题，而不是笼统归因于"逆流"：

\begin{itemize}
\item \textbf{下腹痛伴发热}：警惕盆腔炎，应及时就诊；
\item \textbf{经量骤增、行经期明显延长}：需排查子宫内膜息肉、黏膜下肌瘤、内分泌异常等；
\item \textbf{持续加重的性交痛或痛经}：需评估子宫内膜异位症、腺肌症；
\item \textbf{非经期出血}：需排查宫颈与子宫内膜病变。这些都属于"就医信号"，与经期是否性交无关——但出现时不要拖。
\end{itemize}

\begin{tcolorbox}[colback=pink!5!white,colframe=red!50!black,title=贴心小叮咛]
本节的处理方式，适用于所有"听起来很有道理的传统说法"：把\textbf{理论、证据、推论}三层分开看。理论成立，不代表推论成立；推论成立，也不代表它在每个人身上都必然发生。
\end{tcolorbox}

